% ============================================================================
\section{Introduction}\label{sec:intro}
% ============================================================================

\paragraph{Lossless convexification.}
Powered-descent guidance asks for a thrust history that takes a lander from its current position and velocity
to rest at a target, subject to thrust-magnitude bounds $0<\rho_1\le\norm{T}\le\rho_2$, pointing and state constraints,
and often minimum fuel or minimum landing error. The lower thrust bound makes the control set an annulus, which is
nonconvex. The lossless convexification (LCvx) programme of A\c{c}\i kme\c{s}e and coauthors
\cite{acikmese2005gnc,acikmese2007jgcd,acikmese2011automatica,acikmese2013tcst,blackmore2010jgcd} replaces the annulus
by a convex cone in an augmented space, $\norm{T}\le\Gamma$, $\rho_1\le\Gamma\le\rho_2$, and proves that optimal
solutions of the relaxed second-order cone program (SOCP) satisfy $\norm{T}=\Gamma$ almost everywhere. They are
therefore optimal for the original problem. Extensions cover pointing constraints \cite{carson2011acc,acikmese2013tcst},
state constraints \cite{harris2013cdc,harris2014automatica}, some nonlinear dynamics \cite{blackmore2012scl},
normality-based general theory \cite{harris2021tac}, free final time \cite{kunhippurayil2021automatica} and discrete
time \cite{luo2026automatica}. Surveys and tutorials are \cite{malyuta2021survey,malyuta2022csm}. LCvx-based guidance
has flown \cite{acikmese2013flight,scharf2014aero}. Every one of these results treats the thrust direction as a free
control variable, possibly restricted to a cone about a fixed or prescribed axis. Six-degree-of-freedom (6-DOF)
powered-descent methods use sequential convex programming instead \cite{szmuk2018gnc,szmuk2020jgcd,reynolds2020dq,reynolds2020ptr,elango2025automatica},
and LCvx appears there only as an embedded relaxation of the thrust bounds \cite{malyuta2022csm}. A search of the
literature through November 2026 found no LCvx theorem for a model in which the thrust direction is a
rate-limited state, and no impossibility result for lossless convexification of any kind.

\paragraph{Why 5-DOF and 6-DOF matter.}
On a gimballed (TVC) lander the thrust acts along an engine axis that can deflect by at most a gimbal angle $\delta$ from
the body axis $b$. The body axis itself turns at a finite rate, and on a pure-TVC vehicle it is turned \emph{by} the
gimbal. A 3-DOF LCvx solution can ask for thrust-direction changes the vehicle cannot follow. The lander that
motivated this study is a small throttleable TVC vehicle ($\approx 80$~kg wet, thrust 400--1000~N, $\pm15^\circ$
box gimbal, usable body rate 0.3--0.5~rad/s). Its stored 3-DOF LCvx solutions contain thrust-direction jumps of
$15^\circ$ within one 24~ms control step, plus $180^\circ$ azimuth flips, which would require angular accelerations
around $10^3$~rad/s$^2$. The vehicle can deliver 1--3.5~rad/s$^2$. The question we address is whether \emph{any}
LCvx-type convex program can be lossless once the attitude dynamics are part of the model.

\paragraph{The answer in one paragraph.}
3-DOF LCvx works because the reachable set of a linear system with a measurable control is convex: whatever the
relaxed problem does on average can be realized by fast switching (Lyapunov--Aumann \cite{lyapunov1940,aumann1965}).
In 5-DOF the thrust direction is tied to a body axis that rotates at a finite rate, so it cannot switch. Take a maximal
tilt to the left and a maximal tilt to the right. The average of these two manoeuvres is not flyable. Any convex
program that contains both must also contain their average, provided the physical state is a decision variable and
the boundary data enter affinely. That is exactly the LCvx situation. The rest of the paper makes this precise,
removes the idealizations one by one (state constraints, minimum throttle, terminal attitude, final time, 6-DOF),
quantifies the size of the obstruction, and describes what can still be convexified exactly.

\paragraph{Contributions.}
\begin{enumerate}[leftmargin=*,itemsep=2pt]
\item \emph{Definitions} (\cref{sec:formulation}). We define uniform and state-linear convex relaxations with
affine (or separable-monotone) data entry, and lossless relaxations in a feasibility (detection-sound) form and in
an optimization form. We show that the standard LCvx definition with attainment implies the feasibility form. We also
state precisely what the class excludes: occupation-measure and moment lifts with nonaffine data entry,
per-instance restrictions, and the trivial convex-hull reformulation that every problem admits.
\item \emph{Impossibility} (\cref{sec:impossibility}). Theorem~A: a lossless uniform convex relaxation forces the
feasible boundary-data set to be convex. Theorem~B and an explicit ball bound: the 5-DOF feasible set is nonconvex,
with a gap of metres per second. Theorem~B$^*$: nonconvexity persists with altitude and glide-slope constraints at
every horizon, with upright landing, and with free final time on compact windows. A line-search lemma shows that
searching over $t_f$ cannot rescue LCvx. Theorem~D and Lemma~S give per-instance versions for minimum-landing-error
problems, with a Wirtinger-type inequality that handles coning manoeuvres. Theorem~A-opt and Proposition~T give the
optimization version. Results for $\rho_1=0$ show that the obstruction does not depend on a minimum throttle.
Transfers to kinematic and independent-torque 6-DOF models are given, and pure TVC is treated through a
centre-of-percussion reduction, with its exact scope stated.
\item \emph{Possibility} (\cref{sec:possibility}). LCvx is exact for any fixed attitude profile. For the
gimbal-torque-coupled TVC body the problem with a given attitude is exactly convex, with no relaxation. There is an
exact subclass for the kinematic model. For degree-2 SDP lifts we give stress certificates and a certified-global
coplanar subclass, together with counterexamples to exactness in general.
\item \emph{Numerics} (\cref{sec:numerics}). Gap laws, results beyond the proven regimes, and the size of the
obstruction for a small lander: centimetre-scale for in-air descents, a metre-scale hole in the landing set for
pad-launched hops.
\item \emph{Formalization} (\cref{sec:lean}). Theorem~A, Theorem~B, the state-constrained Theorem~B$^*$, the abstract
Theorem~D and the line-search lemma are machine-checked in Lean~4 with Mathlib \cite{demoura2021lean4,mathlib2020}.
External theorems used on the positive side (the maximum principle and Aumann's theorem) enter only as explicit
hypotheses.
\end{enumerate}

\paragraph{Scope, honestly stated.}
The impossibility concerns the class of convex programs in which the physical state is a decision variable and the
boundary data enter affinely. That covers every LCvx formulation in the literature. It does not cover convex
programs over measures with point-mass data entry, which are exact but infinite-dimensional and per instance, and it
does not say that 5-DOF guidance is intractable. For the in-air descents of the motivating vehicle, the
nonconvexity we can exhibit is real but thin. For the actual pure-TVC actuation with an upright terminal attitude, the
rigorous proof is still being completed (\cref{sec:tvc}).

\paragraph{Notation.}
$\ip{\cdot}{\cdot}$ and $\norm{\cdot}$ are Euclidean, $\Sph$ is the unit sphere, $\ez$ points up, $x_z=\ip{x}{\ez}$ and
$x_h=x-x_z\ez$. For unit vectors, $\angle(x,y)=\arccos\ip xy\in[0,\pi]$. Gravity is $g=-g_0\ez$ unless stated
otherwise. All functions live on $[0,t_f]$ with Lebesgue measure.
